\chapter{Géométrie différentielle}
\minisommairech

Soit $\mathcal{M}$ un espace topologique séparé à base dénombrable d'ouverts, soit $m\in \N$ et soit $r\in \N \cup \{\infty\}$.

\section{Généralités}

\subsubsection{Définitions:}

\begin{defi}[Carte]
Une carte sur $\mathcal{M}$ à valeur dans $\R ^m$ est la donnée d'un couple $(U, \phi)$ où $U$ est un ouvert de $\mathcal{M}$ et $\phi : U \to \phi(U) \in \R^m$ est un homéomorphisme de $U$ sur son image avec $\phi(U)$ ouvert dans $\R^m$.
\end{defi}

\begin{defi}[Atlas]
Un $\mathscr{C}^r$-atlas de $\mathcal{M}$ dans $\R^m$ est un ensemble $\mathscr{A}$ de cartes $(U, \phi)$ à valeur dans $\R^m$ dont les domaines de cartes $U$ recouvrent $\mathcal{M}$ et qui sont $\mathscr{C}^r$-compatibles, c'est à dire telles que:
$$\forall (U, \phi), (V, \psi) \in \mathscr{A}, \quad \psi \circ \phi ^{-1} : \phi(U\cap V) \to \psi(U \cap V) \quad \text{est un }\mathscr{C}^r\text{-difféomorphisme}$$
Deux $\mathscr{C}^r$-atlas sont $\mathscr{C}^r$-compatibles si leur réunion est un $\mathscr{C}^r$-atlas. Cela définit une relation d'équivalence et la réunion de la classe d'équivalence d'un $\mathscr{C}^r$-atlas est l'unique $\mathscr{C}^r$-atlas maximal qui le contient.
\end{defi}

\begin{defi}[Variété différentielle]
Une variété différentielle de classe $\mathscr{C}^r$ (ou lisse si $r = \infty$) de dimension $m$ est un espace topologique séparé et à base dénombrable d'ouverts muni d'un $\mathscr{C}^r$-atlas maximal à valeur dans $\R^m$.
\end{defi}

\begin{defi}
Soient $\mathcal{M}$, $\mathcal{N}$ deux variétés de classe $\mathscr{C}^r$ et $f: \mathcal{M} \to \mathcal{N}$. On dira que $f$ est $\mathscr{C}^r$ (ou lisse si $r = \infty$) si $\forall x\in \mathcal{M}$ il existe une carte $(U, \phi)$ de $\mathcal{M}$ telle que $x\in U$ et une carte $(V, \psi)$ de $\mathcal{N}$ telle que $f(x) \in V$ et telles que $f(U)\subset V$ et si l'application $f$ lue dans les cartes locales 
$$\psi \circ f \circ \phi^{-1}:\phi(U)\to \psi(V)$$
est $\mathscr{C}^r$ en $\phi(x)$. \\
Par composition des applications différentiables, cette définition ne dépend pas du choix des cartes.
\end{defi}

\begin{pro}
Soient $\mathcal{M}$, $\mathcal{N}$ et $\mathcal{P}$ des variétés $\mathscr{C}^r$ et $f:\mathcal{M}\to \mathcal{N}$ et $g:\mathcal{N}\to\mathcal{P}$ des applications $\mathscr{C}^r$. Alors $g \circ f: \mathcal{M} \to \mathcal{P}$ est $\mathscr{C}^r$.
\end{pro}
\begin{preu}
Par composition des applications différentiables
\end{preu}

\subsubsection{Sous-variétés:}

\begin{defi}[Sous-variété]
Soit $\mathcal{M}$ une variété lisse de dimension $m$ et soit $\mathcal{N}$ une partie de $\mathcal{M}$ telle qu'il existe $n\leq m$ tel que pour toute carte $(U, \phi)$ de $\mathcal{M}$ telle que $U \cap \mathcal{N} \neq \emptyset$ on ait
$$\phi : U\cap \mathcal{N} \to \phi(U)\cap \R^n$$
est un homéomorphisme.\\
Dans ce cas on dit que $\mathcal{N}$ est une sous-variété lisse de $\mathcal{M}$ de dimension $n$ et de codimension $m-n$ dans $\mathcal{M}$.
\end{defi}

\begin{pro}
Soit $\mathcal{N}$ une sous variété lisse de $\mathcal{M}$ alors si on note $\{(U_\alpha, \phi_\alpha), \alpha \in A\}$ l'atlas lisse de $\mathcal{M}$ alors l'ensemble $\{(U_\alpha \cap \mathcal{N}, \phi_\alpha), \alpha \in A\}$ est un atlas lisse sur $\mathcal{N}$ qui fait de $\mathcal{N}$ une variété lisse de dimension $n$. \\
De plus, l'inclusion $\mathcal{N} \to \mathcal{M}$ est lisse.
\end{pro}

\subsubsection{Variétés quotients:}

\begin{defi}[Variété quotient]
Soit $\mathcal{M}$ une variété lisse et soit une action lisse libre et propre d'un groupe de Lie $G$ sur $\mathcal{M}$. Alors il existe une unique structure de variété lisse sur $\mathcal{M}/G$ telle que la projection canonique soit une application lisse.
\end{defi}

\section{Fibrés principaux et fibrés vectoriels}

\begin{defi}[Cocycle]
Soit $\mathcal{M}$ une variété lisse, soit $\{(U_\alpha, \phi_\alpha), \alpha \in A\}$ sont atlas lisse et soit $G$ un groupe de Lie. Un cocycle sur $\mathcal{M}$ à valeur dans $G$ est une famille de fonctions lisses $g_{\beta\alpha} : U_\alpha \cap U_\beta \to G$ telle que $\forall \alpha, \beta, \gamma \in A$,
$$g_{\gamma\alpha}(x) = g_{\gamma\beta}(x)g_{\beta\alpha}(x), \quad \forall x \in U_\alpha \cap U_\beta \cap U_\gamma$$ 
\end{defi}

\begin{pro}
Soit $\rho : G \to H$ un morphisme de groupes de Lie et soit un cocycle $g_{\beta\alpha}$ à valeur dans un groupe de Lie $G$, alors $\rho \circ g_{\beta\alpha}$ est un cocycle sur $\mathcal{M}$ à valeur dans $H$. \\
\end{pro}

Soit $\mathcal{M}$ une variété lisse de dimension $m$ et d'atlas $\{(U_\alpha, \phi_\alpha), \alpha \in A\}$ munie d'un cocycle $g_{\beta\alpha}$ à valeur dans un groupe de Lie $G$ de dimension $r$. Notons $V_\alpha = \phi_\alpha(U_\alpha) \subset \R^m$, $V_{\beta\alpha} = \phi_\alpha(U_\alpha \cap U_\beta) \subset \R^m$ et notons $\phi_{\beta\alpha} = \phi_\beta \circ \phi_\alpha^{-1} : V_{\beta\alpha} \to V_{\alpha\beta}$ les fonctions de transitions.

\subsubsection{Fibré principal:}

On pose $H_\alpha = V_\alpha\times G$ et $H_{\beta\alpha} = V_{\beta\alpha}\times G$ et on définit $\gamma_{\beta\alpha} : H_{\beta\alpha} \to H_{\alpha\beta}$ pour tout $\alpha, \beta \in A$ par
$$\gamma_{\beta\alpha}(\phi_\alpha(x), h) = (\phi_\beta(x), g_{\beta\alpha}h), \quad \forall x\in V_{\beta\alpha}, \quad \forall h\in G$$
Soit alors $X = \DUnion_{\alpha \in A} H_\alpha$ muni de la topologie union disjointe et soit la relation binaire sur $X$ définie par
$$(x, g) \in H_{\beta\alpha} \sim (y, h) \in H_{\alpha\beta} \iff (y,h) = \gamma_{\beta\alpha}(x,g)$$
C'est une relation d'équivalence sur $X$ puisque la famille des $\gamma_{\beta\alpha}$ vérifie:
\be
	\item $\gamma_{\alpha\alpha} = \id$ sur $H_\alpha$
	\item $\gamma_{\beta\alpha}^{-1} = \gamma_{\alpha\beta}$ sur $H_{\alpha\beta}$
	\item $\gamma_{\delta\beta} \circ \gamma_{\beta\alpha} = \gamma_{\delta\alpha}$ sur $H_{\beta\alpha} \cap \gamma_{\beta\alpha}^{-1}(H_{\delta\beta}) \subset H_{\delta\alpha}$ \\
\ee

Soit alors $P = X/\sim$ muni de la topologie quotient. Soit alors $\{(\mathcal{O}_\beta, \theta_\beta), \beta \in B\}$ l'atlas de $G$, l'application $V_\alpha \times \theta_\beta(\mathcal{O}_\beta) \subset \R^m \times \R^r \to P$ qui pour $x\in \mathcal{M}$ et $h\in G$ envoie $(\phi_\alpha(x), \theta_\beta(h))$ sur sa classe d'équivalence dans $P$ est un homéomorphisme sur son image. Notons $p_{\alpha, \beta} : P_{\alpha, \beta} \to V_\alpha \times \theta_\beta(\mathcal{O}_\beta)$ l'application réciproque. \\

Ainsi, l'ensemble $\{(P_{\alpha, \beta}, p_{\alpha, \beta}), \alpha\in A, \beta\in B\}$ définit alors un atlas lisse sur $P$ qui est alors une variété lisse de dimension $m+r$. \\

De plus, les applications $\pi_\alpha: H_\alpha \to V_\alpha$ de projection sur la première coordonnée induisent une surjection lisse $\pi : P \to \mathcal{M}$, lisse par définition de l'atlas de $P$. Et on dispose de $\mathscr{C}^\infty$-difféomorphismes $\Phi_\alpha : \pi^{-1}(U_\alpha) \to H_\alpha$ appelés trivialisations locales et telles qu'on a le diagramme commutatif:
$$\begin{matrix}
P \supset \pi^{-1}(U_\alpha) & \overset{\Phi_\alpha}{\longrightarrow} & H_\alpha \\ \\
\pi\big\downarrow &  & \big\downarrow\pi_\alpha \\ \\
\mathcal{M} \supset U_\alpha & \overset{\phi_\alpha}{\longrightarrow} & V_\alpha 
\end{matrix}
$$
Enfin, on dispose d'une action à droite lisse de $G$ sur chaque $H_\alpha$ définie pour $g\in G$ et $(\phi_\alpha(x), h)\in H_\alpha$ par 
$$(\phi_\alpha(x), h).g = (\phi_\alpha(x), hg) \in H_\alpha$$
Et pour $(\phi_\alpha(x), h) \in H_{\beta\alpha}$ et $g\in G$ on a 
$$[\gamma_{\beta\alpha}(\phi_\alpha(x), h)].g = (\phi_\beta(x), g_{\beta\alpha}hg) = \gamma_{\beta\alpha}[(\phi_\alpha(x), h).g]$$ 
Ainsi cela définit une action à droite lisse de $G$ sur $P$ qui laisse stable chaque fibre $P_x$ au dessus de $x\in \mathcal{M}$. 

\begin{defi}[Fibré principal]
On appelle fibré principal sur $\mathcal{M}$ et de groupe structural $G$ l'application $\pi: P \to \mathcal{M}$ qui vérifie:
\be 
	\item Il existe des difféomorphismes $\Phi_\alpha : \pi^{-1}(U_\alpha) \to V_\alpha \times G$ appelés trivialisations locales tels que $\pi_\alpha \circ \Phi_\alpha = \phi_\alpha \circ \pi$. 
	\item $G$ agit à droite de façon lisse sur $P$ tel que chaque fibre $P_x = \pi^{-1}(x)$ soit stable par $G$ et tel que $G$ agisse transitivement sur chaque fibre $P_x$.
\ee
On a alors $\dim P = \dim \mathcal{M} + \dim G$ et on note $G \hookrightarrow P \overset{\pi}{\to} \mathcal{M}$.
\end{defi}

\begin{defi}[Morphisme de fibrés principaux]
Soient $G \hookrightarrow P \overset{\pi}{\to} \mathcal{M}$ et $G \hookrightarrow Q \overset{\rho}{\to} \mathcal{N}$ deux fibrés principaux. Un morphisme de fibré principaux est la donnée d'une application lisse $f:\mathcal{M}\to \mathcal{N}$ et d'une application lisse $F: P \to Q$ stable par l'action de $G$ telle que le diagramme suivant commute:
$$\begin{matrix}
P  & \overset{F}{\longrightarrow} & Q \\ \\
\pi\big\downarrow &  & \big\downarrow\rho \\ \\
\mathcal{M} & \overset{f}{\longrightarrow} & \mathcal{N} 
\end{matrix}
$$
On dit que $F$ recouvre $f$. \\

On dit que $G \hookrightarrow P \overset{\pi}{\to} \mathcal{M}$ et $G \hookrightarrow Q \overset{\rho}{\to} \mathcal{N}$ sont équivalent s'il existe un isomorphisme $P\to Q$ de fibrés principaux qui recouvre $\id_\mathcal{M}$.
\end{defi}

\begin{defi}[Fibré trivial]
$\pi : \mathcal{M}\times G \to \mathcal{M}$ qui est la projection par rapport à la première composante s'appelle le fibré principal trivial explicite. \\
Un fibré principal sur $\mathcal{M}$ de groupe structural $G$ est dit trivial s'il est équivalent au fibré principal trivial explicite. 
\end{defi}

\begin{defi}[Section globale, section locale]
Soit $G \hookrightarrow P \overset{\pi}{\to} \mathcal{M}$ un fibré principal sur $\mathcal{M}$. Une section globale est une application lisse $s:\mathcal{M} \to P$ telle que $\pi \circ s = \id_\mathcal{M}$ ou telle que $\forall x\in \mathcal{M}, s(x) \in P_x$. \\
Pour tout ouvert $U\subset \mathcal{M}$, on appelle section locale sur $U$ toute section de $\pi : \pi^{-1}(U) \to U$. \\
On note $\Gamma(P)$ ou $\Gamma(\tau)$ l'ensemble des sections de $\tau: E \to \mathcal{M}$
\end{defi}

\begin{pro}
Un fibré principal admet une section globale si et seulement si il et trivial. De plus, l'ensemble des sections globales est en bijection avec l'ensemble des trivialisations.
\end{pro}

\subsubsection{Fibrés vectoriels:}

Soit $F$ un $\R$-espace vectoriel de dimension finie et soit $\rho : G \to \GL(F)$ une représentation de $G$. On pose $W_\alpha = V_\alpha\times F$ et $W_{\beta\alpha} = V_{\beta\alpha}\times F$ et on définit $\psi_{\beta\alpha} : W_{\beta\alpha} \to W_{\alpha\beta}$ par
$$\psi_{\beta\alpha}(\phi_\alpha(x), v) = (\phi_\beta(x), \rho(g_{\beta\alpha})v), \quad \forall x\in V_{\beta\alpha}, \quad \forall v\in F$$

Soit alors $Y = \DUnion_{\alpha \in A} W_\alpha$ muni de la topologie union disjointe et soit la relation binaire sur $Y$ définie par
$$(x,u)\in W_{\beta\alpha} \sim (y, v)\in W_{\alpha\beta} \iff (y, v) = \psi_{\beta\alpha}(x, u)$$
C'est une relation d'équivalence sur $X$ puisque la famille des $\psi_{\beta\alpha}$ vérifie:
\be
	\item $\psi_{\alpha\alpha} = \id$ sur $W_\alpha$
	\item $\psi_{\beta\alpha}^{-1} = \psi_{\alpha\beta}$ sur $W_{\alpha\beta}$
	\item $\psi_{\delta\beta} \circ \psi_{\beta\alpha} = \psi_{\delta\alpha}$ sur $W_{\beta\alpha} \cap \psi_{\beta\alpha}^{-1}(W_{\delta\beta}) \subset W_{\delta\alpha}$ \\
\ee

Soit alors $E = X/\sim$ muni de la topologie quotient. L'application $W_\alpha \to E$ qui envoie les éléments de $W_\alpha$ sur leur classe d'équivalence dans $E$ est un homéomorphisme sur son image. Notons, $e_\alpha: E_\alpha \to W_\alpha$ l'application réciproque. \\

Ainsi, l'ensemble $\{(E_\alpha, e_\alpha), \alpha \in A\}$ définit alors un atlas lisse sur $E$ dont les fonctions de transition sont exactement les applications $\psi_{\beta\alpha}$ ce qui fait de $E$ une variété lisse de dimension $m+k$. \\

De plus, les applications $\tau_\alpha: W_\alpha \to V_\alpha$ de projection sur la première coordonnée induisent une surjection lisse $\tau: E \to \mathcal{M}$. Notons alors $E_x = \tau^{-1}(x)$ la fibre au dessus de $x$. Soit $(U_\alpha, \phi_\alpha)$ une carte autour de $x$ dans $\mathcal{M}$ et soit $(E_\alpha, e_\alpha)$ la carte associée dans $E$. Alors, $e_\alpha : \{x\}\times E_x \to \{\phi_\alpha(x)\} \times F$ est une bijection. On muni $E_x$ d'une structure de $\R$-espace vectoriel tel que cette bijection devienne une application linéaire. Cette définition est indépendante du choix de la carte par construction des fonctions de transition $\psi_{\beta\alpha}$. \\

\begin{defi}[Fibré vectoriel]
On appelle fibré vectoriel sur $\mathcal{M}$ de fibre modèle $F$ l'application $\tau: E \to \mathcal{M}$ telle que pour tout $x\in \mathcal{M}$ la fibre $E_x$ est un $\R$-espace vectoriel isomorphe à $F$. On a alors $\dim E = \dim \mathcal{M} + \dim F$. \\
C'est la fibration vectorielle associée au fibré principal $G \hookrightarrow P \overset{\pi}{\to} \mathcal{M}$ via la représentation $\rho : G \to \GL(F)$, notée parfois $P\times_\rho F$.
\end{defi}

\begin{defi}[Sous-fibré vectoriel]
Soit $\tau: E \to \mathcal{M}$ un fibré vectoriel de fibre modèle $V$. Un sous-fibré vectoriel de $\tau$ est $\tau_{|F}$ où $F$ est une partie de $E$ telle que 
\be
	\item $E_x \cap F$ est un sous-espace vectoriel de $E_x$ pour tout x$\in \mathcal{M}$
	\item Pour tout $x \in \mathcal{M}$, il existe un voisinage ouvert $U$ de $x$, une trivialisation locale $\varphi: \tau^{-1}(U) \to U\times V$ et un sous-espace vectoriel $W$ de $V$ tel que $\phi(F \cap \tau^{-1}(U)) = U \times W$ \\
\ee 
\end{defi}

\begin{defi}[Produit fibré]
Soient $\tau : E = \to \mathcal{M}$ et $\sigma : F = \to \mathcal{M}$ des fibrés vectoriels, on pose 
$$E\diamond F = \{(u,v) \in E\times F | \tau(u) = \sigma(v)\}$$
Alors $p: E\diamond F \to \mathcal{M}$ est un fibré vectoriel appelé produit fibré de $E$ et $F$.
\end{defi}

Soit $G \hookrightarrow P \overset{\pi}{\to} \mathcal{M}$ un fibré principal et soient $\tau : E = P\times_\rho V \to \mathcal{M}$ et $\tau' : F = P\times_{\rho'} W  \to \mathcal{M}$ deux fibrés vectoriels. \\

\begin{defi}[Fibré vectoriel dual]
On appelle dual du fibré vectoriel $E$ et on note $E*$ le fibré vectoriel $P \times_{\rho*} (V*)$ donné par la représentation duale $\rho*$.
\end{defi}

\begin{defi}[Somme directe de fibrés vectoriels]
On appelle somme directe des fibré vectoriels $E$ et $F$ et on note $E\oplus F$ le fibré vectoriel $P \times_{\rho \oplus \rho'} (V\oplus W)$ donné par la représentation somme directe $\rho \oplus \rho'$.
\end{defi}

\begin{defi}[Produit tensoriel de fibrés vectoriels]
On appelle produit tensoriel des fibré vectoriels $E$ et $F$ et on note $E\otimes F$ le fibré vectoriel $P \times_{\rho \otimes \rho'} (V\otimes W)$ donné par la représentation produit tensoriel $\rho \otimes \rho'$.
\end{defi}

\begin{defi}[Section]
Soit $\tau: E \to \mathcal{M}$ un fibré vectoriel sur $\mathcal{M}$. Une section est une application $s:\mathcal{M} \to E$ telle que $\pi \circ s = \id_\mathcal{M}$ ou de façon équivalente telle que $\forall x\in \mathcal{M}, s(x) \in E_x$. \\
On note $\Gamma(E)$ ou $\Gamma(\tau)$ l'espace vectoriel des sections de $\tau: E \to \mathcal{M}$. C'est aussi un $\mathscr{C}^\infty(\mathcal{M})$-module.
\end{defi}

\subsubsection{Fibré tangent, fibré cotangent et fibré tensoriel:}

Notons $t_{\beta\alpha} = d\phi_{\beta\alpha} \circ \phi_\alpha : U_\alpha \cap U_\beta \to \GL(\R^m)$ le cocycle canoniquement donné par l'atlas de $\mathcal{M}$.

\begin{defi}[Fibré tangent]
Le fibré tangent de $\mathcal{M}$ noté $\tau_\mathcal{M} : T\mathcal{M} \to \mathcal{M}$ est donné par le cocycle $t_{\beta\alpha}$ et la représentation fondamentale de $\GL(\R^m)$. \\
On appelle champs de vecteur une section du fibré $\tau_\mathcal{M}$.
\end{defi}

\begin{defi}[Fibré cotangent]
Le fibré cotangent de $\mathcal{M}$ noté $\tau^*_\mathcal{M} : T^*\mathcal{M} \to \mathcal{M}$ est donné par le cocycle $t_{\beta\alpha}$ et la représentation duale de $\GL(\R^m)$. 
\end{defi}

\begin{defi}[Fibré tensoriel]
Soient $r,s\in \N$, alors le fibré tensoriel d'ordre $(r,s)$ ($r$ covariant et $s$ contravariant) de $\mathcal{M}$, noté $\tau^{(r,s)}_\mathcal{M} : T^{(r,s)}\mathcal{M} \to \mathcal{M}$ est donné par le cocycle $t_{\beta\alpha}$ et la représentation tensorielle $\GL(\R^m) \to \GL(\otimes^r(\R^m)^* \otimes^s\R^m)$ de la représentation fondamentale et de sa représentation duale. \\
On appelle champs tensoriel une section du fibré $\tau^{(r,s)}_\mathcal{M}$.
\end{defi}

\section{Formes différentielles}

\begin{defi}[Forme différentielle]
Notons $\Lambda^k T^*\mathcal{M}$ donné par la représentation $\GL(\R^m) \to \GL(\Lambda^k(\R^m)^*), u\mapsto \Lambda^k u^*$ et $\Lambda T^*\mathcal{M} = \oplus_{k\geq 0} \Lambda^k T^*\mathcal{M}$. \\
Une $k$-forme differentielle sur $\mathcal{M}$ est une section du fibré $\Lambda^k T^*\mathcal{M}$. On note alors $\Omega^k\mathcal{M} = \Gamma(\Lambda^k T^*\mathcal{M})$ et $\Omega \mathcal{M} = \Gamma(\Lambda T^*\mathcal{M})$. \\

On dispose de l'application $\wedge : \Omega^p\mathcal{M} \times \Omega^q\mathcal{M} \to \Omega^{p+q}\mathcal{M}$ appelé produit extérieur qui est induit par le produit extérieur sur les espaces $\Lambda^k(\R^m)^*$. Ainsi, on muni l'espace $\Omega\mathcal{M}$ d'une structure d'algèbre. \\

Soit $\tau: E \to \mathcal{M}$ un fibré vectoriel de fibre modèle $V$, une $k$-forme différentielle à valeur dans $V$ est une section du fibré $\Lambda^k T^*\mathcal{M}\otimes E$. On note alors $\Omega^k(\mathcal{M}, E) = \Gamma(\Lambda^k T^*\mathcal{M}\otimes E)$ et $\Omega(\mathcal{M}, E) = \Gamma(\Lambda T^*\mathcal{M}\otimes E)$. \\

On dispose également d'un produit extérieur induit par celui sur les applications multilinéaires alternée $\wedge:\Lambda^p(\R^m)^* \otimes E \times \Lambda^q(\R^m)^* \otimes F \to \Lambda^{p+q}(\R^m)^* \otimes E \otimes F$  
$$\wedge: \Omega^p(\mathcal{M}, E) \times \Omega^q(\mathcal{M}, F) \to \Omega^{p+q}(\mathcal{M}, E\otimes F)$$
\end{defi}

\begin{rem}
L'identification entre $\Lambda^k(\R^m)^* \otimes E$ et les application $k$-linéaires alternées de $(\R^m)^k$ à valeur dans $E$ permet d'identifer les $k$-formes différentielles sur $\mathcal{M}$ à valeur dans $E$ avec la donnée sur chaque $T_x\mathcal{M}$ d'une application $k$-linéaire alternée de $T_x\mathcal{M}^k$ à valeur dans $E$. Et le produit extérieur ainsi définit correspond au produit extérieur d'applications multilinéaires alternées sur chaque $T_x\mathcal{M}$. 
\end{rem}

\begin{defi}
Soient $\mathcal{M}, \mathcal{N}$ deux variétés lisse, $f:\mathcal{N} \to \mathcal{M}$ une application lisse et $\eta \in \Omega^k\mathcal{M}$. Alors $f^*\eta = \eta (Tf, ..., Tf) : T^{\diamond k}\mathcal{N} \to \R$ est une $k$-forme différentielle sur $\mathcal{N}$ appelée la forme tirée en arrière par $f$ de $\eta$. \\
\end{defi}

\begin{defi}[Différentielle extérieure]
Soit $(U, \phi)$ une carte de $\mathcal{M}$ et soit $x\in U$ et notons $\phi = (x^1, ..., x^m)$ et notons $(\partial_1, ..., \partial_m)$ la base de $T_x\mathcal{M}$ image reciproque de la base canonique de $\R^m$ par $T_x\phi$ et notons $(dx^1, ..., dx^m)$ la base duale de $(\partial_1, ..., \partial_m)$ dans $(T_x\mathcal{M})^*$. Pour $f\in \Omega^0\mathcal{M} = \mathscr{C}^\infty(\mathcal{M})$ on définit l'application linéaire $d : \Omega^0\mathcal{M} \to \Omega^1\mathcal{M}$ par 
$$df = \frac{\partial f \circ \phi^{-1}}{\partial x^1}dx^1 +...+ \frac{\partial f \circ \phi^{-1}}{\partial x^m}dx^m$$
En particulier on a $d(x^k) = dx^k$. \\
Puis pour $\omega \in \Omega^k\mathcal{M}$ une $k$-forme sur $\mathcal{M}$ s'écrivant $\omega = f dx^{i_1} \wedge ... \wedge dx^{i_k}$ alors on définit l'application linéaire $d: \Omega^k\mathcal{M} \to \Omega^{k+1}\mathcal{M}$ par 
$$d\omega = df \wedge dx^{i_1} \wedge ... \wedge dx^{i_k}$$
On appelle les application $d: \Omega^k\mathcal{M} \to \Omega^{k+1}\mathcal{M}$ la différentielle extérieure sur $\mathcal{M}$.
\end{defi}

\begin{pro}
La dérivée extérieure $d: \Omega^k\mathcal{M} \to \Omega^{k+1}\mathcal{M}$ vérifie les propriétés suivantes :
\be
	\item $d$ est linéaire
	\item $d^2 = 0$ 
	\item Pour $\omega_1 \in \Omega^p\mathcal{M}$ et $\omega_2 \in \Omega^q\mathcal{M}$, $\quad d(\omega_1 \wedge \omega_2) = d\omega_1 \wedge \omega_2 + (-1)^p\omega_1 \wedge d\omega_2$
	\item Pour $f\in \Omega^0\mathcal{M}$ et pour tout champs de vecteurs $X$ sur $\mathcal{M}$ on a $df(X) = fX$
	\item Pour $f:\mathcal{N} \to \mathcal{M}$ lisse et pour $\omega \in \Omega^k\mathcal{M}$ on a $f^*(d\omega) = d(f^*\omega)$
	\item Pour tout champs de vecteurs $X,Y$ sur $\mathcal{M}$ et pour toute $1$-forme $\omega \in \Omega^1\mathcal{M}$ on a $d\omega(X,Y) = X(\omega(Y)) - Y(\omega(X)) - \omega([X, Y])$
\ee
\end{pro}

\begin{defi}
On dit qu'une $k$-forme $\omega$ est fermée si $d\omega = 0$. \\
On dit qu'elle est exacte s'il existe une $k-1$-forme $\eta$ telle que $\omega = d\eta$. \\
Puisque $d^2 = 0$ on a que toute forme exacte est fermée.
\end{defi}

\begin{pro}[Lemme de Poincaré]
Localement toute forme différentielle fermée est exacte.
\end{pro}

\begin{defi}[Cohomologie de Rham]
Soit $\mathcal{M}$ une variété lisse, pour $k\in \N^*$ on définit la $k$-ème cohomologie de Rham l'espace quotient 
$$H_{dr}^k\mathcal{M} = \frac{\Ker d_{k+1}}{\Ima d_k}$$
où $d_k = d: \Omega^k\mathcal{M} \to \Omega^{k+1}\mathcal{M}$.
\end{defi}

\begin{exe}
Soit $U$ est un ouvert de $\R^3$, on identifie $(\R^3)^*$ et $\R^3$. Alors pour $f \in \Omega^0U = \mathscr{C}^\infty(U)$ on a:
$$df = \frac{\partial f}{\partial x} dx + \frac{\partial f}{\partial y} dy + \frac{\partial f}{\partial z} dz = \grad f$$
Soit $\omega = f_x dx + f_y dy + f_z dz \in \Omega^1U$, alors on a
$$d\omega = \left(\frac{\partial f_z}{\partial y} - \frac{\partial f_y}{\partial z}\right)dy\wedge dz + \left(\frac{\partial f_x}{\partial z} - \frac{\partial f_z}{\partial x}\right)dz\wedge dx + \left(\frac{\partial f_y}{\partial x} - \frac{\partial f_x}{\partial y}\right)dx\wedge dy$$
or $\star (dy\wedge dz) = dx$, $\quad\star (dz\wedge dx) = dy$ et $\quad\star (dx\wedge dy) = dz$ et on a 
$$\star d\omega = \rot f$$ 
Soit $\omega = f_x dy\wedge dz + f_y dz\wedge dx + f_z dx \wedge dy \in \Omega^2U$, alors on a 
$$d\omega = \left(\frac{\partial f_x}{\partial x} + \frac{\partial f_y}{\partial y} + \frac{\partial f_z}{\partial z}\right) dx\wedge dy\wedge dz$$
or $\star (dx\wedge dy\wedge dz) = 1$ et
$$\star d\omega = \Div f$$
Finalement on a 
$$\begin{matrix}
\Omega^0 & \overset{d}{\longrightarrow} & \Omega^1 & \overset{d}{\longrightarrow} & \Omega^2 & \overset{d}{\longrightarrow} & \Omega^3\\
\Big\updownarrow &  & \Big\updownarrow &  & \Big\updownarrow\star &  & \Big\updownarrow\star\\
\R & \overset{\grad}{\longrightarrow} & \R^3 & \overset{\rot}{\longrightarrow} & \R^3 & \overset{\Div}{\longrightarrow} & \R
\end{matrix}
$$
En utilisant $d^2 = 0$ on a immediatement 
$$\rot(\grad f) = 0 $$
et 
$$\Div(\rot f) = 0$$
\end{exe}

\section{Flot local d'un champs de vecteurs}

\begin{defi}[Flot local d'un champs de vecteurs]
Soit $X$ un champs de vecteurs sur $\mathcal{M}$. Un flot local de $X$ est une application $\phi : \Omega \to \mathcal{M}$ lisse notée $(t,x) \to \phi_t(x)$ où $\Omega$ est un voisinage ouvert de $\{0\}\times \mathcal{M}$ dans $\R \times \mathcal{M}$ tel que pour tout $x\in \mathcal{M}$, $\Omega \cap \R \times \{x\}$ soit connexe, et qui vérifie les deux propriétés suivantes: 
$$\frac{d\phi_t(x)}{dt}_{|t=s} = X(\phi_s(x)) \quad \forall (s, x)\in \Omega$$
$$\phi_0(x) = x \quad \forall x\in \mathcal{M}$$
Il est dit complet si $\Omega = \R \times \mathcal{M}$. \\
Il est dit maximal s'il n'existe pas de prolongement strict. Autrement dit tel que s'il existe $\phi' : \Omega' \to \mathcal{M}$ un flot local de $X$ tel que $\Omega \subset \Omega'$ et $\phi'|_{\Omega} = \phi$ alors $\Omega = \Omega'$ et $\phi = \phi'$. Dans ce cas il est alors unique.
\end{defi}

\begin{defi}[Courbe integrale]
Soit $X$ un champs de vecteurs sur $\mathcal{M}$. Une courbe integrale de $X$ est une application $\gamma : I \to \mathcal{M}$ telle que $\dot\gamma = X \circ \gamma$ sur $I$. \\
Elle est dite maximale s'il n'existe pas de prolongement strict. Autrement dit tel que s'il existe une autre courbe integrale $c: J \to \mathcal{M}$ tel que $I \subset J$ et $c|_I = \gamma$ alors $I = J$ et $c = \gamma$.
\end{defi}

\begin{pro}
Notons $\supp(X) = \overline{\{x\in \mathcal{M} | X(x) \neq 0\}}$. Alors 
\be
	\item S'il existe $\varepsilon > 0$ tel que $]-2\varepsilon; 2\varepsilon[\times \mathcal{M} \subset \Omega$ alors $\phi_t^X$ est complet.
	\item Si $\supp(X)$ est compact alors $\phi_t^X$ est complet.
\ee
\end{pro}

\begin{pro}
Un champs de vecteurs $X$ sur $\mathcal{M}$ admet un unique flot maximal $\phi_t^X : \Omega \to \mathcal{M}$ qui vérifie de plus les propriétés suivantes:
\be
	\item $\forall x_0\in \mathcal{M}$, l'unique courbe integrale maximale de $X$ passant par $x_0$ au temps $t=0$ est la courbe $I_{x_0} = \{t\in \R | (t, x_0) \in \Omega\} \to \mathcal{M}$ définie par
	$$t \mapsto \phi_t^X(x_0)$$

	\item $\forall (s, x)\in \Omega$, on a $(t, \phi_s^X(x)) \in \Omega \iff (t+s, x) \in \Omega$ et dans ce cas
	$$\phi_t^X \circ \phi_s^X(x) = \phi_{t+s}^X(x)$$

	\item $\forall t\in \R$ l'application $\phi_t^X : U_t = \{x\in \mathcal{M} | (t, x) \in \Omega \} \to \mathcal{M}$ est un $\mathscr{C}^\infty$-difféomorphisme local. Si de plus $X$ est complet alors $(\phi_t^X)_{t\in\R}$ est un groupe à un paramètre de $\mathscr{C}^\infty$-difféomorphismes de $\mathcal{M}$. 

	\item $\forall (t,x) \in \Omega$, le flot local préserve le champs de vecteurs, c'est à dire 
	$$T_x\phi_t^X(X(x)) = X(\phi_t^X(x))$$
\ee
\end{pro}

\section{Dérivation de Lie}

\begin{defi}[Dérivation]
Une dérivation de $\mathcal{M}$ est une application $\R$-linéaire $\delta : \mathscr{C}^\infty(\mathcal{M}) \to \mathscr{C}^\infty(\mathcal{M})$ telle que pour tout $f, g \in \mathscr{C}^\infty(\mathcal{M})$ on ait
$$\delta(fg) = f\delta(g) + g\delta(f)$$
On note $\mathscr{D}(\mathcal{M})$ l'ensemble des dérivations sur $\mathcal{M}$. C'est un $\R$-espace vectoriel et c'est aussi un $\mathscr{C}^\infty(\mathcal{M})$-module.
\end{defi}

\begin{defi}[Dérivation de Lie]
Soit $X$ un champs de vecteurs sur $\mathcal{M}$ alors l'application $\mathscr{L}_X: \mathscr{C}^\infty(\mathcal{M}) \to \mathscr{C}^\infty(\mathcal{M})$, appelée la dérivation de Lie associée à $X$ et notée parfois $X(f) = \mathscr{L}_X(f)$, est définie par 
$$\mathscr{L}_X(f): x\to T_xf(X(x))$$
Le théorème de dérivation des fonctions composées montre que si $\phi$ est un flot local sur $X$ alors
$$\mathscr{L}_X(f) = \frac{d}{dt}_{|t = 0}f\circ\phi_t$$
\end{defi}

\begin{pro}
L'application $X \in \Gamma(T\mathcal{M}) \to \mathscr{L}_X \in \mathscr{D}(\mathcal{M})$ est un isomorphisme de $\R$-espaces vectoriels ainsi que de $\mathscr{C}^\infty(\mathcal{M})$-modules.
\end{pro}

\begin{defi}[Crochet de champs de vecteurs]
Soient $\delta, \delta'$ deux dérivations sur $\mathcal{M}$ alors
$$[\delta, \delta'] = \delta \circ \delta' - \delta' \circ \delta$$
est une dérivation sur $\mathcal{M}$. \\

Si $X, Y$ sont deux champs de vecteurs sur $\mathcal{M}$ on note $[X, Y]$ le champs de vecteurs tel que 
$$\mathscr{L}_{[X, Y]} = [\mathscr{L}_X, \mathscr{L}_Y]$$
\end{defi}

\begin{pro}
Le crochet de champs de vecteurs $[X, Y]$ muni $\Gamma(T\mathcal{M})$ d'une structure d'algèbre de Lie telle que l'application $\mathscr{L}$ soit un isomorphisme d'algèbres de Lie.
\end{pro}

\section{Théorème de Frobenius}

\subsubsection{Fibré grassmannien:}

\begin{defi}[Grasmannienne]
Soient $0\leq k \leq n$. On définie la grasmannienne des $k$-espaces de $\R^n$ par 
$$G(n,k) = \{ W \subset \R^n | \text{ tel que W est un sous-espace vectoriel de dimension k}\}$$ 
On peut voir $G(n,k) = \GL(\R^n) / \Sigma$ où $\Sigma = \{f\in \GL(\R^n) | f(\R^k) = \R^k\}$. Ainsi, $G(n,k)$ est muni d'une structure de variété lisse. \\
\end{defi} 

Soit $\mathcal{M}$ une variété lisse et soit $t_{\beta\alpha} : U_\alpha \cap U_\beta \to \GL(\R^m)$ le cocycle de son fibré tangent où $\{(U_\alpha, \phi_\alpha), \alpha \in A\}$ désigne l'atlas lisse de $\mathcal{M}$. On définit alors les fonctions de transitions 
$$\psi_{\beta\alpha} : \phi_\alpha(U_\alpha\cap U_\beta) \times G(n, k) \to \phi_\beta(U_\alpha \cap U_\beta) \times G(n,k)$$
donné par 
$$\psi_{\beta\alpha}(\phi_\alpha(x), V) = (\phi_\beta(x) , t_{\beta\alpha}(x)V)$$
pour tout $x\in U_\alpha \cap U_\beta \subset \mathcal{M}$ et pour tout $V\in G(n,k)$. \\

\begin{defi}[Fibré grassmannien]
Cela définit un fibré $\pi : G(T\mathcal{M}, k) \to \mathcal{M}$ appelé le fibré grasmannien de $k$-espace sur $\mathcal{M}$. \\
Une section de $\pi : G(T\mathcal{M}, k) \to \mathcal{M}$ est appelée une distribution lisse de $k$-espaces sur $\mathcal{M}$.
\end{defi}

\begin{defi}
Soit $\mathcal{M}$ une variété lisse, soit $W$ une distribution lisse de $k$-espaces sur $\mathcal{M}$ et soit $x_0\in \mathcal{M}$. On dit qu'une sous variété $\mathcal{N}$ de $\mathcal{M}$ de dimension $k$ intègre $W$ avec condition initiale $x_0$ si $x_0 \in \mathcal{N}$ et $T_x\mathcal{N} = W(x)$ pour tout $x\in \mathcal{N}$. \\
On dit que $W$ est intégrable s'il existe une telle sous-variété $\mathcal{M}$ pour tout $x_0\in \mathcal{M}$.
\end{defi}

\subsubsection{Théorème de Frobénius:}

\begin{defi}[Condition d'intégrabilité de Frobenius]
Soit $W$ une distribution de $k$-espaces sur une variété lisse $\mathcal{M}$. On dit que $W$ satifait la condition d'intégrabilité de Frobenius dans un ouvert $U\in \mathcal{M}$ si pour tout champs de vecteurs $X,Y$ définit sur $U$ tels que $X(x) \in W(x)$ et $Y(x) = W(x)$ pour tout $x\in U$ on a $[X, Y](x) \in W(x)$ pour tout $x\in \mathcal{M}$.
\end{defi}

\begin{pro}[Théorème d'intégrabilité de Frobenius]
Soit $W$ une distribution de $k$-espaces sur une variété lisse $\mathcal{M}$. Alors $W$ est intégrable si et seulement si elle vérifie la condition d'intégrabilité de Frobenius.
\end{pro}

\section{Connexions}

Soit $G \hookrightarrow P \overset{\pi}{\to} \mathcal{M}$ un fibré principal de groupe structural $G$ sur $\mathcal{M}$. On pose pour $p\in P$
$$i_p : G \to P, \quad g \mapsto p.g$$
Et comme $i_p(e) = p$, on dispose de 
$$T_e i_p: \mathfrak{g} \to T_pP$$

\begin{defi}[Vecteur vertical]
On dit qu'un vecteur tangent de $T_pP$ est vertical s'il est dans l'image de $T_e i_p$. On note 
$$\Verti(T_pP) = \Ima(T_e i_p) = T_e i_p(\mathfrak{g})$$
De plus $p\in P \mapsto \Verti(T_pP)$ est une distribution de $(\dim G)$-espaces sur $P$. 
\end{defi}

\begin{pro}
On a la suite exacte courte d'espaces vectoriels
$$0 \to \mathfrak{g} \overset{T_e i_p}{\longrightarrow} T_p P \overset{T_p \pi}{\longrightarrow} T_x\mathcal{M} \to 0$$ 
Ainsi on a $\dim(\Verti(T_pP)) = \dim G$. \\
\end{pro}

\begin{defi}
Pour $X\in \mathfrak{g}$ on pose $X^\sharp(p) = T_e i_p(X) \in \Verti(T_pP)$. Puis on définit 
$$X^\sharp: P \to TP, \quad p \mapsto X^\sharp(p)$$
On dit que $X^\sharp$ est le champs de vecteur fondamental sur $P$ associé à $X$. \\
De plus on a l'application $.^\sharp : \mathfrak{g} \to \Gamma(TP)$ est un morphisme d'algèbres de Lie.
\end{defi}

\subsubsection{Connection d'Ehresmann:}

\begin{defi}[Connexion d'Ehresmann]
Notons $R_g : p\in P \mapsto p.g \in P$ et soit $G \hookrightarrow P \overset{\pi}{\to} \mathcal{M}$ un fibré principal. Une connexion sur $P$ est la donnée $p\in P \mapsto \Hor(T_pP)$ d'une distribution lisse de $m$-espaces sur $P$ telle que:
\be
	\item $\Verti(T_pP) \oplus \Hor(T_pP) = T_pP$
	\item $T_pR_g(\Hor(T_pP)) = \Hor(T_{p.g}P)$ pour tout $g\in G$ et tout $p\in P$
\ee 
Ou de façon équivalente, c'est la donnée d'une 1-forme différentielle $\omega$ sur $\mathcal{M}$ à valeur dans $\mathfrak{g}$ telle que
$$R_g^*\omega  = \Ad_{g^{-1}} \circ \omega \quad \forall g\in G$$
ou autrement dit si elle vérifie la condition 
$$\omega(X^\sharp) = X \quad \forall X \in \mathfrak{g}$$
\end{defi}
\begin{preu}
Soit $p \in P \mapsto \Hor(T_pP)$ une connexion sur $P$. Cela fournit une application linéaire $v_p : T_pP \to \Verti(T_pP)$ qui est la projection sur la première composante de la somme directe. Alors pour chaque $p\in P$ on dispose d'une application linéaire 
$$\omega_p : T_pP \overset{v_p}{\longrightarrow} \Verti(T_pP) \overset{(T_e i_p)^{-1}}{\longrightarrow} \mathfrak{g}$$
Qui fourni une 1-forme différentielle sur $P$ à valeur dans $\mathfrak{g}$
$$\omega : TP \to \mathfrak{g}$$
\end{preu}

\begin{pro}
Notons $x = \pi(p)$ alors d'après la suite exacte courte, l'application $T\pi$ fourni un isomorphisme d'espaces vectoriels
$$\Hor(T_pP) \to T_x\mathcal{M}$$
\end{pro}

\begin{pro}
On a 
\be
	\item $R_g \circ i_p = i_{p.g} \circ \ad_{g^{-1}}$
	\item $T_pR_g(\Verti(T_pP)) = \Verti(T_{p.g}P)$
\ee
\end{pro}
\begin{preu}
On a $R_g \circ i_p(h) = p.hg = p.gg^{-1}hg = i_{p.g}(ad_{g_{-1}}(h))$. D'où (i).\\
(ii) s'obtient en considérant les applications tangentes dans (i).
\end{preu}

\begin{defi}
Soient $\mathcal{M}, \mathcal{N}$ deux variétés lisse, $V$ un $\R$-espace vectoriel, $f:\mathcal{N} \to \mathcal{M}$ une application lisse et $\eta \in \Omega^1(\mathcal{M}, V)$. Alors $f^*\eta = \eta \circ Tf : T\mathcal{N} \to V$ est une $1$-forme différentielle sur $\mathcal{N}$ à valeur dans $V$ appelée la forme tirée en arrière par $f$ de $\eta$. \\
\end{defi}

\begin{pro}
Soient $G \hookrightarrow P \overset{\pi}{\to} \mathcal{M}$ et $G \hookrightarrow G \overset{\rho}{\to} \mathcal{N}$ des fibrés principaux de groupe structural $G$ et soit $F: P\to Q$ un morphisme de fibré principaux recouvrant l'application $f: \mathcal{M} \to \mathcal{N}$ et soit $\omega$ une connexion d'Ehresmann sur $Q$. Alors $F^*\omega$ est une connexion d'Ehresmann sur $P$.
\end{pro}

\begin{pro}
Soit $\pi : P \to \mathcal{M}$ un fibré principal de groupe structural $G$ muni d'une connexion d'Ehresmann $\omega$, soit $\gamma : [0,1] \to \mathcal{M}$ une courbe lisse et soit $p_0\in \pi^{-1}(\gamma(0))$. Alors il existe une unique courbe lisse $\tilde{\gamma_{p_0}} : [0,1] \to P$ telle que 
\be
	\item $\tilde{\gamma_{p_0}}(0) = p_0$
	\item $\pi \circ \tilde{\gamma_{p_0}} = \gamma$
	\item $\tilde{\gamma_{p_0}}'(t) \in \Hor(T_{\tilde{\gamma_{p_0}}(t)}P)$ pour tout $t \in [0,1]$
\ee
\end{pro}

\begin{pro}
Soit $G \hookrightarrow P \overset{\pi}{\to} \mathcal{M}$ un fibré principal muni d'une connexion d'Ehresmann $\omega$, soit $\gamma : [0,1] \to \mathcal{M}$ une courbe lisse. On définit une application 
$$\Par_t^\gamma : \pi^{-1}(\gamma(0)) \to \pi^{-1}(\gamma(t)), \quad p_0 \mapsto \tilde{\gamma_{p_0}}(t)$$
\end{pro}

\subsubsection{Forme de Maurer-Cartan:}

\begin{defi}[Forme de Maurer-Cartan]
Soit $G$ un groupe de Lie. On appelle forme de Maurer-Cartan sur $\mathfrak{g}$, la 1-forme à valeur dans $\mathfrak{g}$ notée $\Omega : TG \to \mathfrak{g}$, et définit par 
$$\Omega_g = T_gL_{g^{-1}} : T_gG \to \mathfrak{g}$$
pour tout $g\in G$.
\end{defi}

\begin{pro}
Soit $G$ un groupe de Lie. Alors sa forme de Maurer-Cartan est l'unique connexion d'Ehresmann sur le fibré principal $G \to \{x\}$.
\end{pro}

\begin{pro}
La forme de Maurer-Cartan d'un groupe $G$ fournit une trivialisation globale $\tilde{\Omega} : TG \to G\times \mathfrak{g}$ de $TG$ définie par
$$\tilde{\Omega}(p) = (\tau_G(p), \Omega(p))$$
\end{pro}

\subsubsection{Courbure:} 

\begin{defi}[Courbure]
Soit $G \hookrightarrow P \overset{\pi}{\to} \mathcal{M}$ un fibré principal muni d'une connexion d'Ehresmann $\omega$. La courbure de $\omega$ est l'application $D\omega : T^{\diamond 2}P \to \mathfrak{g}$ définie par 
$$D\omega(u,v) = d\omega(u^H, v^H) \quad \forall (u,v)\in T^{\diamond 2}P$$
\end{defi}

\begin{pro}
Soit $G \hookrightarrow P \overset{\pi}{\to} \mathcal{M}$ un fibré principal muni d'une connexion d'Ehresmann $\omega$. Alors la distribution associée $p\to \Hor(T_pP)$ est intégrable si et seulement si $D\omega = 0$.
\end{pro}

\begin{pro}[Equation de structure de Cartan]
Soit $G \hookrightarrow P \overset{\pi}{\to} \mathcal{M}$ un fibré principal muni d'une connexion d'Ehresmann $\omega$. Alors on a 
$$D\omega = d\omega + \frac{1}{2}\omega\wedge\omega$$
\end{pro}

\begin{pro}[Identité de Bianchi]
Soit $G \hookrightarrow P \overset{\pi}{\to} \mathcal{M}$ un fibré principal muni d'une connexion d'Ehresmann $\omega$. Notons $\Theta = D\omega$ alors
$$D\Theta = 0$$
\end{pro}